\documentclass[11pt]{article}
\usepackage[a4paper,margin=28mm]{geometry}
\usepackage{amsmath,amssymb,amsthm,mathtools,hyperref}
\newcommand{\ddc}{dd^c}
\newcommand{\dc}{d^c}
\newcommand{\MA}{\operatorname{MA}}
\newcommand{\PSH}{\operatorname{PSH}}
\newcommand{\PP}{\mathbb P}
\newtheorem{theorem}{Theorem}
\newtheorem{lemma}[theorem]{Lemma}
\theoremstyle{remark}\newtheorem{remark}[theorem]{Remark}
\title{Exponential equations at zero individual threshold\newline
for a constructed concentration family}
\author{}\date{}
\begin{document}\maketitle
The general classification in S03 remains unresolved here.
\begin{abstract}
We construct a strictly positive integrable probability density on a
compact Hermitian product whose individual Monge--Amp\`ere threshold
is zero, but whose exponential equation has a genuine unbounded
subsolution for every positive exponent. The density is a rigorously
chosen member of the specified concentration family and is independent
of the exponent. An explicit radial density-growth criterion supplies
the subsolutions, with exactly computed finite mass. The supplied
subsolution principle then gives greatest finite-mass exact solutions,
which are necessarily unbounded. A prior elliptic threefold supplies
the required smooth mass collapse and concrete examples in dimensions
four and five. No assertion is made for every concentration-family
choice or for necessity of the sufficient criterion.
\end{abstract}

\section{Conventions, background and product geometry}
Use $\ddc=i\partial\bar\partial/(2\pi)$ and
$\dc=i(\bar\partial-\partial)/(4\pi)$, with ordinary wedge powers.
All quasi-psh potentials have their genuine canonical representatives.
Extended measures are non-pluripolar products with smooth-background
corrections. For a positive Hermitian form $h$ and reference probability
$dV$, let $c_h(g)$ be the unique scalar of the normalized
bounded-density equation, and set
\[
 b_h(f)=\lim_{j\to\infty}c_h(\min(f,j)),\qquad
 v_-(h)=\inf_{v\in\PSH(X,h)\cap L^\infty}\int_X\MA_h(v).
\]
For nonzero nonnegative bounded data, existence of a continuous
potential and uniqueness of this scalar follow from
\cite[Theorem 0.1]{KN} and \cite[Corollary 0.2]{Ng}.
Their $\ddc$ is twice ours, so their potentials are multiplied by two;
the smooth positive reference-volume ratio is absorbed into the density.
Neither theorem is applied to the final unbounded $L^1$ density.
We use the supplied maximal-subsolution characterization of $b_h$.
In particular, $\MA_h^{\rm ext}(v)\geq af\,dV$ implies $a\leq b_h(f)$,
and bounded data satisfy $b_h=c_h$.

We also use exactly the supplied S03 principle: a genuine exponential
subsolution $s$ with
$\MA_h^{\rm ext}(s)\geq e^{\lambda s}f\,dV$, $\lambda>0$, gives a
greatest genuine exact solution by decreasing bounded-density
approximation. This principle is background; we prove the required
subsolutions and their zero-threshold nonvacuity below.

Let $(X_0,h_0)$ be connected compact Hermitian of dimension $m$,
and put $Y=\PP^k$, $k\geq1$, with normalized Fubini--Study form
$\kappa$, $\int_Y\kappa^k=1$. On $X=X_0\times Y$ write
\[
 N=m+k,\quad C=\binom Nm,\quad V_0=\int_{X_0}h_0^m,
 \quad H=h_0+\kappa,
 \quad\mu=\frac{h_0^m}{V_0}\wedge\kappa^k.
\]
For any smooth positive probability $dV$ on $X$, put $R=\mu/dV$.
It is smooth with positive upper and lower bounds. Factor forms and
functions are pulled back when appropriate. Write $F=f/R$, so that
$f\,dV=F\mu$.

\section{A finite-mass radial barrier}
Fix $0<\alpha<1$ and $o=[1:0:\cdots:0]$. Put
\[
 \psi=\log\frac{\sum_{i=1}^k|Z_i|^2}{\sum_{i=0}^k|Z_i|^2},
 \qquad L=1-\psi,\qquad u=1-L^\alpha,
\]
with the canonical value $u(o)=-\infty$. For $1\leq L<\infty$ set
\begin{align*}
 t&=e^{1-L},& a&=\alpha L^{\alpha-1},\\
 Q&=t+(1-t)a,&
 D&=t(1-a)+(1-t)\alpha(1-\alpha)L^{\alpha-2},\\
 q_\alpha&=Q^{k-1}D e^{k(L-1)},&
 W&=kQ^{k-1}D.
\end{align*}
Here $q_\alpha$ is a density function on $Y\setminus\{o\}$;
its extension over $Z_0=0$ has value $1-\alpha$.

\begin{lemma}\label{lem:barrier}
The potential $u$ is genuine $\kappa$-psh and
\[
 \MA_\kappa^{\rm ext}(u)=q_\alpha\kappa^k,\qquad
 \int_Yq_\alpha\kappa^k=1,\qquad q_\alpha\geq(1-\alpha)^k.
\]
Its product pullback has actual extended mass $CV_0$. Moreover,
as $L\to\infty$,
\[
 W(L)\sim k\alpha^k(1-\alpha)L^{-1-k(1-\alpha)},\qquad
 e^{\tau(L^\alpha-1)}W(L)\longrightarrow\infty\quad(\tau>0).
\]
\end{lemma}
\begin{proof}
For $0<\varepsilon\leq1$ replace $\psi$ by the smooth function
\[
 \psi_\varepsilon=
 \log\frac{\sum_{i=1}^k|Z_i|^2+\varepsilon|Z_0|^2}
               {\sum_{i=0}^k|Z_i|^2}.
\]
It is nonpositive and $\kappa$-psh. For
$\chi(s)=1-(1-s)^\alpha$ one has $0<\chi'\leq\alpha$ and
$\chi''>0$, so
\[
 \kappa+\ddc\chi(\psi_\varepsilon)
 =(1-\chi')\kappa+\chi'(\kappa+\ddc\psi_\varepsilon)
   +\chi''d\psi_\varepsilon\wedge\dc\psi_\varepsilon
 \geq(1-\alpha)\kappa.
\]
These functions decrease to $u$. Its growth
$O(1+|\log r|^\alpha)$, $r=\sum_{i=1}^k|z_i|^2$ near $o$,
is locally integrable. Thus the decreasing psh-plus-smooth limit is
nontrivial and canonical, with its stated pole.

In the affine chart $Z_0=1$, $L=1-\log(r/(1+r))$ and a local
primitive of $\kappa+\ddc u$ is
\[
 \Phi(r)=\log(1+r)+1-[1-\log(r/(1+r))]^\alpha,
 \qquad r\Phi'=Q.
\]
The tangential and radial Hessian eigenvalues are $Q/r$ and $Q'$.
Polar integration, with its constant fixed by normalized
Fubini--Study mass, gives
\[
 \int_{a<r<b}(\ddc\Phi)^k=Q(b)^k-Q(a)^k.
\]
Since $-dQ/dL=D$, the radial mass density is $W(L)dL$.
The Fubini--Study radial probability is $ke^{-k(L-1)}dL$;
division gives exactly the displayed $q_\alpha$.
The flux limits are $Q(0)=0$, $Q(\infty)=1$ in the $r$ coordinate.
The function is smooth along the infinity hyperplane. On every
finite superlevel, its truncation agrees locally with smooth $u$;
the defining non-pluripolar exhaustion omits its pole and cutoff
masses. Hence the actual mass is one. The lower form bound proves
the density bound. Factor dimensions give the same finite-level
identity on the product, followed by the actual non-pluripolar limit:
\[
 \MA_H^{\rm ext}(u)=C h_0^m\wedge q_\alpha\kappa^k=CV_0q_\alpha\mu.
\]
Finally $Q\sim\alpha L^{\alpha-1}$ and
$D\sim\alpha(1-\alpha)L^{\alpha-2}$ give the asymptotic for $W$.
The logarithm of the last expression in the lemma is
$\tau L^\alpha-[1+k(1-\alpha)]\log L+O(1)$, tending to infinity.
\end{proof}

For $\tau>0$ write $B_\tau=q_\alpha e^{\tau(L^\alpha-1)}$.
In particular $B_\tau\geq b_0=(1-\alpha)^k$.
Its mass on every punctured pole neighborhood is infinite, even
though $q_\alpha\kappa^k$ has mass one.

\begin{theorem}\label{th:criterion}
Let $f\geq0$ be nonzero and integrable, and fix $\lambda>0$.
If, for some finite $A>0$,
\[
 F(x,y)\leq A B_\lambda(y)\quad\text{almost everywhere},
\]
then the genuine function
\[
 s_\lambda=u+\lambda^{-1}\log(CV_0/A)
\]
is an exponential subsolution with actual mass $CV_0$ and exact
coefficient $\lambda$. The supplied principle gives a greatest exact
solution. The same holds for every $\lambda>0$ if $F$ is locally
essentially bounded off $X_0\times\{o\}$ and
\[
 \mathop{\rm ess\,limsup}_{y\to o}\mathop{\rm ess\,sup}_{x}
    \frac{\log^+(F(x,y)/q_\alpha(y))}{L(y)^\alpha}=0.
\]
\end{theorem}
\begin{proof}
The exact product measure in Lemma~\ref{lem:barrier} gives
\[
 e^{\lambda s_\lambda}f\,dV
 \leq(CV_0/A)e^{\lambda u}Aq_\alpha e^{-\lambda u}\mu
 =\MA_H^{\rm ext}(s_\lambda).
\]
Both measures give zero mass to the polar fiber, and adding the
constant changes neither admissibility nor mass. For the second
claim, the limit condition bounds $F/q_\alpha$ by
$e^{(\lambda/2)L^\alpha}$ near $o$. The remaining compact region
has bounded ratio, so increasing a finite $A$ gives the first claim.
The greatest solution has finite mass because it is bounded above
on compact $X$ and $f$ is integrable.
\end{proof}

The first criterion can also be checked from the explicit growth bound
\[
 F\leq A'e^{k(L-1)+\lambda L^\alpha}L^{-1-k(1-\alpha)}
\]
for all sufficiently large $L$, uniformly essentially in $x$ and
the projective spheres, with local essential boundedness elsewhere.

\section{A zero-threshold concentration family for every exponent}
Assume that $X_0$ has smooth positive probability densities $F_j$
relative to $dV_0=h_0^m/V_0$ satisfying
\begin{equation}\label{eq:base}
 c_{h_0}(F_j)\leq 2^{-2j}/C.
\end{equation}
\begin{theorem}\label{th:construction}
For every sequence satisfying \eqref{eq:base}, one can choose smooth
strictly positive probability densities $g_j$ on $Y$ such that
\[
 f_j=RF_jg_j,\qquad f=\sum_{j\geq1}2^{-j}f_j
\]
are the specified full-product concentration family:
$c_H(f_j)\leq2^{-2j}$, $\int_Xf\,dV=1$, and $b_H(f)=0$.
The density $f$ is finite everywhere, smooth off the polar fiber and
bounded below by $\tfrac12\min_Xf_1>0$. This single $f$, independent
of $\lambda$, has genuine exponential subsolutions and greatest
finite-mass exact solutions for every $\lambda>0$. Every exact
solution is unbounded.
\end{theorem}
\begin{proof}
Put $w_j=2^{-j}$ and $M_j=\max F_j$. Choose
$0<\varepsilon_j\leq1/2$ so that
\[
 w_j\varepsilon_jM_j\leq b_0 2^{-j-2},\qquad
 w_j\varepsilon_j\|F_j\|_{C^j}\leq2^{-j},
\]
using increasing fixed smooth coordinate seminorms. Then
$E(x)=\sum_jw_j\varepsilon_jF_j(x)$ converges smoothly and
$0<E\leq b_0/4$.

Choose a smooth $0\leq\rho\leq1$ supported in $[0,1]$, positive
inside and flat at the endpoints. Inductively choose
$A_j\geq2$, $A_j>A_{j-1}+2$, with
\[
 Z_j=\int_{A_j}^{A_j+1}\rho(L-A_j)e^{(L^\alpha-1)/j}W(L)\,dL
      \geq2w_jM_j.
\]
For each fixed $j$ the integrand apart from $\rho$ tends to infinity
uniformly on intervals $[A,A+1]$ by Lemma~\ref{lem:barrier}.
Thus the choice is possible with arbitrarily large $A_j$.
Set, on that annulus and zero elsewhere,
\[
 b_j=\rho(L-A_j)B_{1/j}/Z_j,\qquad
 g_j=\varepsilon_j+(1-\varepsilon_j)b_j.
\]
Flatness makes $b_j$ smooth globally, vanishing near both $o$ and
the infinity hyperplane. Its integral is one. The $g_j$ are smooth
strictly positive probabilities and their disjoint annuli satisfy
\begin{equation}\label{eq:bump}
 w_jM_j(1-\varepsilon_j)b_j\leq B_{1/j}/2.
\end{equation}

We check the component scalar directly. For a radial $g=g_j$, put
\[
 G(r)=\int_{|z|^2<r}g\kappa^k,\quad Q_g=G^{1/k},\quad
 \phi'(r)=\frac{Q_g-r/(1+r)}r.
\]
The tangential and radial eigenvalues of the local primitive are
$Q_g/r>0$ and $Q_g'>0$, and the flux identity gives
$(\kappa+\ddc\phi)^k=g\kappa^k$. At small $r$, the bump is absent:
$G=\varepsilon_j(r/(1+r))^k$ and
$\phi=(\varepsilon_j^{1/k}-1)\log(1+r)+\mathrm{constant}$.
At large $r$,
\[
 G=1-\varepsilon_j+\varepsilon_j(r/(1+r))^k.
\]
In $v=1/r$, $Q_g$ is analytic near zero with value one and
$d\phi/dv=-(Q_g-1/(1+v))/v$ is analytic there. Since
$v=|Z_0|^2/\sum_{i\geq1}|Z_i|^2$ is smooth near $Z_0=0$,
$\phi$ extends smoothly. The metric is strictly positive there,
with tangential part $\kappa$ and transverse coefficient
$\varepsilon_j>0$; at $o$ it is $\varepsilon_j^{1/k}\kappa$.
This includes $k=1$.

Let $v_j$ solve the bounded base equation above with scalar
$c_{h_0}(F_j)$. The genuine bounded product potential
$U_j=v_j+\phi_j$ has
\[
 \MA_H(U_j)=C\MA_{h_0}(v_j)\wedge g_j\kappa^k
          =C c_{h_0}(F_j)f_j\,dV.
\]
The local bounded-product expansion, including smooth-background
corrections, has only this term by factor dimensions. Scalar
uniqueness and \eqref{eq:base} therefore give
$c_H(f_j)=C c_{h_0}(F_j)\leq2^{-2j}$.

Tonelli gives $\int f=\sum_jw_j=1$. Off the polar fiber, only finitely
many annuli meet any compact set, so
$F=E+\sum_jw_j(1-\varepsilon_j)F_jb_j$ is smooth. On the fiber
the actual pointwise sum equals the finite $E$.
Also $f\geq w_1f_1$, proving the positive lower bound.
If $a>0$ were an admissible subsolution coefficient for $f$, then
for every $j$ maximality would give
\[
 aw_j\leq b_H(f_j)=c_H(f_j)\leq2^{-2j},
 \qquad a\leq2^{-j},
\]
which is impossible. Thus $b_H(f)=0$.

Fix $\lambda>0$ and choose $J$ with $1/j\leq\lambda$ for $j\geq J$.
Since $u\leq0$, $B_{1/j}\leq B_\lambda$ for these indices.
Disjointness and \eqref{eq:bump} bound their total contribution
pointwise by $B_\lambda/2$. The baseline is at most $B_\lambda/4$.
The finitely many remaining bumps have finite constant
\[
 K_\lambda=\max_Y\sum_{j<J}w_j(1-\varepsilon_j)M_jb_j,
 \qquad D_\lambda=3/4+K_\lambda/b_0.
\]
Hence $F\leq D_\lambda B_\lambda$.
Theorem~\ref{th:criterion} gives the explicit subsolution
$s_\lambda=u+\lambda^{-1}\log(CV_0/D_\lambda)$ and a greatest
exact solution. The subsolution mass is $CV_0$; no such numerical
mass is asserted for the greatest solution. Its exact mass is finite
and positive, since its supremum is finite, $f\in L^1$ is strictly
positive, and a genuine nontrivial quasi-psh function is finite
almost everywhere. A bounded exponential subsolution would give
a positive constant coefficient for $f$, contradicting $b_H(f)=0$.
Every exact solution is therefore unbounded.
\end{proof}

The construction also satisfies the second criterion of
Theorem~\ref{th:criterion}: on the $j$th annulus,
$F/q_\alpha\leq1/4+\tfrac12 e^{(L^\alpha-1)/j}$, and elsewhere the
ratio is at most $1/4$. As the pole is approached through annuli,
$j\to\infty$, giving the stated uniform logarithmic limit.

\section{Concrete zero-lower-mass examples and scope}
The prior supplied factor in \cite[Appendix]{Product} is
\[
 X_0=\{[z:w]\in\PP^3:z,w\in\mathbb C^2\setminus\{0\}\}
       /\langle[z:w]\mapsto[z:e^{2\pi}w]\rangle,
 \quad h_0=\alpha_0+\beta_0+dt\wedge\dc t,
\]
where $t=\log(\|w\|^2/\|z\|^2)$ and the two base forms are
normalized. Its complete prior proof establishes compactness,
positivity, $V_0=24\pi$ and actual smooth $v_\ell$ with
$h_0+\ddc v_\ell>0$ and
$\delta_\ell=\int(h_0+\ddc v_\ell)^3\to0$.
Choose a subsequence with $\delta_{\ell_j}\leq2^{-2j}/C$ and put
\[
 F_j=\frac{(h_0+\ddc v_{\ell_j})^3/dV_0}{\delta_{\ell_j}}.
\]
These are smooth positive probabilities, and their actual equations
give $c_{h_0}(F_j)=\delta_{\ell_j}$ by scalar uniqueness.
Thus Theorem~\ref{th:construction} has a concrete instance.
For bounded base potentials $v$, the same local product expansion gives
$\int_X\MA_H(v)=C\int_{X_0}\MA_{h_0}(v)$, proving $v_-(H)=0$.
With $k=1$ this is a compact fourfold and the explicit subsolutions
have mass $96\pi$; with $k=2$ the mass is $240\pi$.
These are masses of the constructed subsolutions, not of the greatest
exact solutions.

\begin{remark}
The quantifier is every base sequence satisfying \eqref{eq:base},
with projective components chosen by the displayed construction.
It does not cover every previously chosen full-product concentration
family. The fixed $f$ solves the zero-threshold solvability problem
for this choice and every exponent. Necessity, arbitrary Hermitian
geometries and a nonsolvable zero-threshold example remain unproved.
No literature priority or discovery of the prior factor is claimed.
\end{remark}
\begin{thebibliography}{9}
\bibitem{KN} S. Ko\l odziej and N. C. Nguyen,
\emph{Weak solutions to the complex Monge--Amp\`ere equation on
Hermitian manifolds}, \href{https://arxiv.org/abs/1312.5491v2}{arXiv:1312.5491v2}, Theorem 0.1.
\bibitem{Ng} N. C. Nguyen,
\emph{The complex Monge--Amp\`ere type equation on compact Hermitian
manifolds and applications},
\href{https://arxiv.org/abs/1501.00891v1}{arXiv:1501.00891v1}, Corollary 0.2.
\bibitem{Product} \emph{Exact equations for $L^1$ densities on products
with zero lower mass}, checked standalone source (2026), Appendix;
\href{https://github.com/MingchenXia/non-kahler-pluripotential-theory/blob/7a244bc4938c5b918f7d906c7201f9385e4188f4/research/papers/P07-zero-lower-mass-product-equations.tex}{immutable source and complete prior factor proof}.
\end{thebibliography}
\end{document}
